% !PRINTED 321
\weberpage{321}

296. \textit{Disposition der Aufgabe.} Die beiden gegebenen Formen seien
\[
f = a_x^2, \quad \varphi = \alpha_x^3;
\]
das volle System von $f$ besteht aus den zwei Formen
\[
f, \quad A_{ff},
\]
das von $\varphi$ aus den vier Formen
\weq{2}{\varphi, \quad \Delta, \quad Q, \quad R = A_{\Delta\Delta}.}
Durch Ueberschiebung vollständiger Systeme entstehen wieder vollständige Systeme. Das volle simultane System von $f$ und $\varphi$ erhalten wir also, wenn wir Producte
\[
U = f^\varrho \cdot A_{ff}^\sigma
\]
mit Producten
\[
V = \varphi^\varkappa \cdot \Delta^\lambda \cdot Q^\mu \cdot A_{\Delta\Delta}^\nu
\]
überschieben, d. h. mit einander multipliciren und auf alle möglichen Arten zur Faltung bringen. Da bei der Faltung Invarianten nicht betheiligt sind, und da ferner $Q$ Functionaldeterminante von $\varphi$ und $\Delta$ ist, so können wir in diesen Producten von Factoren $A_{ff}$ und $A_{\Delta\Delta}$ absehen und höhere Potenzen von $Q$ nach der Relation
\[
Q^2 = -\frac{1}{2}\{\Delta^3 + A_{\Delta\Delta}\varphi^2\}
\]
durch $\varphi$ und $\Delta$ ersetzt denken. Unsere Aufgabe reducirt sich daher darauf:

\glqq Alle jene nicht zerfallenden Formen zu finden, die aus den Faltungen
\[
(U, \; V)^\alpha = (f^\varrho, \; \varphi^\varkappa \cdot \Delta^\lambda \cdot Q)^\alpha
\]
hervorgehen, wenn $\varrho, \varkappa, \lambda, \alpha$ alle möglichen Werthe annehmen.\grqq{}

Wir gehen hiebei, wie früher, schrittweise vor, indem wir zunächst solche Producte $V$ mit $f^\varrho$ falten, die nur eine Form $\varphi^\mu$, oder $\Delta^\mu$, oder $Q$ enthalten. Erst in zweiter Linie werden Faltungen mit Producten $V$ untersucht, die mehr als eine der drei Formen zu Factoren haben.

297. \textit{Erste Gruppe von Ueberschiebungen.} Wir haben demnach zunächst zu untersuchen: Zu welchen Formen geben die Ueberschiebungen
\[
(f^\varrho, \; \Delta^\mu)^\alpha, \quad (f^\varrho, \; \varphi^\mu)^\alpha, \quad (f^\varrho, \; Q)^\alpha
\]
Veranlassung? Aus dem ersten Operationsymbol ergeben sich nur zwei Formen:
\[
(f, \; \Delta) \text{ und } (f, \; \Delta)^2 = A_{f\Delta},
\]
%CONT>>


% ===== pdf idx 336, printed 322.0 =====
% !PRINTED 322
\weberpage{322}
%<<CONT
wie aus der Theorie zweier quadratischer Formen bekannt ist. Das zweite Operationssymbol führt zunächst zu folgenden sechs Formen:
\[
(f, \; \varphi), \quad (f, \; \varphi)^2, \quad (f^2, \; \varphi)^3, \quad (f^2, \; \varphi^2)^4, \quad (f^3, \; \varphi^2)^5, \quad (f^3, \; \varphi^2)^6.
\]
Von diesen sind aber die vierte und fünfte Form reducibel. Denn
\[
(f^2, \; \varphi^2)^4 \text{ enthält das zerfallende Glied } (f, \; \varphi)^2 \cdot (f, \; \varphi)^2
\]
und
\[
(f^3, \; \varphi^2)^5 \quad \text{\textquotedbl} \quad \text{\textquotedbl} \quad \text{\textquotedbl} \quad \text{\textquotedbl} \quad (f, \; \varphi)^2 \cdot (f^2, \; \varphi)^3.
\]
Die vier anderen Formen weisen dagegen keine zerfallenden Glieder auf. Höhere Potenzen von $f$ und $\varphi$ brauchen nicht mehr überschoben zu werden. Denn jede solche Ueberschiebung $(f^\varrho, \; \varphi^\mu)^\alpha$ würde a priori ein Glied besitzen, welches zerfällt, nämlich das Glied
\[
(f^3, \; \varphi^2)^6 \cdot (f^{\varrho-3}, \; \varphi^{\mu-2})^{\alpha-6}.
\]
Das dritte Operationssymbol endlich liefert überhaupt nur drei Formen
\[
(Q, \; f), \quad (Q, \; f)^2, \quad (Q, \; f^2)^3,
\]
von denen die erste als Functionaldeterminante einer Functionaldeterminante reducibel ist, während die beiden anderen Formen in das System aufgenommen werden müssen.

298. \textit{Zweite Gruppe von Ueberschiebungen.} Wir überschieben nun Producte $V$ über $f^\varrho$ von der Form
\[
V = \Delta\varphi, \quad \Delta Q, \quad \Delta\varphi Q, \quad \varphi Q.
\]
Allein da die drei ersten Producte den Factor $\Delta$ besitzen, so wird jede der Ueberschiebungen
\[
(f^\varrho, \; \Delta\varphi)^\alpha, \quad (f^\varrho, \; \Delta Q)^\alpha, \quad (f^\varrho, \; \Delta\varphi Q)^\alpha
\]
ein Glied enthalten, welches eine Potenz von $A_{f\Delta}$ zum Factor hat, also zerfällt. Es erübrigt daher nur noch, die Ueberschiebungen
\[
(f^\varrho, \; \varphi Q)^\alpha
\]
zu untersuchen. Hiebei muss $\alpha$ a priori grösser als 3 und demnach $\varrho$ mindestens gleich 2 sein. Wir haben somit zunächst die drei Formen:
\[
(f^2, \; \varphi Q)^4, \quad (f^3, \; \varphi Q)^5, \quad (f^3, \; \varphi Q)^6.
\]
Nur die letzte derselben, eine Invariante, enthält kein zerfallendes Glied; die erste dagegen enthält das Glied $(f, \; \varphi)^2 \cdot (f, \; Q)^2$, die zweite das Glied $(f, \; \varphi)^2 \cdot (f^2, \; Q)^3$. Weitere Ueberschiebungen zu untersuchen ist überflüssig; sie besitzen stets Glieder, welche die Invariante $(f^3, \; \varphi Q)^6$ zum Factor haben.

299. \textit{Zusammenfassung der Resultate.} Das simultane vollständige System von $f$ und $\varphi$ besteht demnach aus folgenden 15 Formen:
%CONT>>


% ===== pdf idx 337, printed 323.0 =====
% !PRINTED 323
\weberpage{323}
%<<CONT
\begin{tabular}{@{}ll@{}}
Drei cubische Covarianten: & $\varphi, \; Q, \; (f, \; \varphi) = \vartheta$ \\
Drei quadratische \quad ,, & $f, \; \Delta, \; (f, \; \Delta) = \Theta$ \\
Vier lineare \quad ,, & $(f, \; \varphi)^2, \; (f^2, \; \varphi)^3, \; (f, \; Q)^2, \; (f^2, \; Q)^3$ \\
Fünf Invarianten: & $A_{ff}, \; A_{\Delta\Delta}, \; A_{f\Delta}, \; (f^3, \; \varphi^2)^6, \; (f^3, \; \varphi \cdot Q)^6.$ \\
\end{tabular}

Zwischen den fünf Invarianten muss eine Relation bestehen. Denn die beiden Formen $f$ und $\varphi$ besitzen nur sieben Constante, von denen drei durch lineare Transformation zum Verschwinden gebracht werden können: demnach können nur vier von einander unabhängig sein.

300. \textit{Einführung einfacherer Symbole für die Formen des Systemes.} Für die symbolische Rechnung ist es wiederum vortheilhaft, jene Formen des Systemes, die noch nicht durch ein einfaches symbolisches Product dargestellt sind, durch eines ihrer Glieder zu ersetzen. Solche Formen sind die beiden linearen Covarianten $(f, \; Q)^2$ und $(f^2, \; Q)^3$ und die beiden Invarianten $(f^3, \; \varphi^2)^6$, $(f^3, \; \varphi \cdot Q)^6$. Setzen wir nämlich
\[
(f, \; \varphi)^2 = (a\alpha)^2 \alpha_x = p
\]
\[
(f^2, \; \varphi)^3 = (a\alpha)^2 (b\alpha) b_x = q = (f, \; p),
\]
so lassen sich an Stelle der beiden andern linearen Covarianten $(f, \; Q)^2$ und $(f^2, \; Q)^3$ die Formen:
\weq{1}{(p, \; \Delta) = (p\Delta)\Delta_x = r}
resp.:
\weq{2}{(\Theta, \; p) = (\Theta p)\Theta_x = s}
und an Stelle der beiden erwähnten Invarianten die Formen
\weq{3}{(f, \; p^2)^2 = (ap)^2 = F}
resp.
\weq{4}{(\Theta, \; p)^2 = (sp) = -M}
einführen.

Denn es ist:
\[
(f, \; Q)^2 = (a_x^2, \; (\alpha\Delta)\Delta_x \alpha_x^2)^2
\]
und $(\alpha\Delta)(a\alpha)^2\Delta_x = (p, \; \Delta) = r$ ist ein Glied dieser Ueberschiebung. Es ist ferner:
\[
(f^2, \; Q)^3 = (a_x^2 b_x^2, \; (\alpha\Delta)\Delta_x \alpha_x^2)^3.
\]
Greifen wir auch hier das Glied heraus, das durch Faltung der Factoren $a_x^2$ mit $\alpha_x^2$, $b_x$ mit $\Delta_x$ entsteht, so können wir dasselbe durch:
\[
(a\alpha)^2(\alpha\Delta)(b\Delta)b_x = (p\Delta)(b\Delta)b_x
\]
darstellen. Das symbolische Product $(p\Delta)(b\Delta)b_x$ ist aber ein Glied der Ueberschiebung $(p, \; (b\Delta)b_x\Delta_x) = (p, \; \Theta)$, und demnach lässt sich in der That $(f^2, \; \Theta)^3$ durch $(\Theta, \; p) = -(p, \; \Theta)$ ersetzen.


% ===== pdf idx 338, printed 324.0 =====
% !PRINTED 324
\weberpage{324}
Des weitern besitzt die Invariante $(f^3, \; \varphi^2)^6$ das Glied
\[
(a\alpha)^2(b\beta)^2(c\alpha)(c\beta) = (f, \; p^2)^2 = (ap)^2 = F
\]
und endlich die Invariante $(f^3, \; \varphi \cdot Q)^6 = (a_x^2 b_x^2 c_x^2, \; \alpha_x^3 \cdot (\beta\Delta)\Delta_x \beta_x^2)^6$ das Glied:
\[
(\beta\Delta) \cdot (a\alpha)^2(b\alpha)(b\Delta)(c\beta)^2 = -((b\Delta)b_x\Delta_x, \; p^2)^2 = -(\Theta, \; p)^2.
\]
Die Tabelle des simultanen Formensystemes wird demnach nach Einführung dieser neuen Symbole:
\[
1) \quad \varphi = a_x^3, \qquad Q = (\alpha\Delta)a_x^2\Delta_x, \qquad \vartheta = (a\alpha)a_x\alpha_x^2;
\]
\[
2) \quad f = a_x^2, \qquad \Delta = (\alpha\beta)^2\alpha_x\beta_x, \qquad \Theta = (a\Delta)a_x\Delta_x;
\]
\[
3) \quad p = (a\alpha)^2\alpha_x, \qquad q = (ap)a_x, \qquad r = (p\Delta)\Delta_x,
\]
\[
\qquad s = (\Theta p)\Theta_x;
\]
\[
4) \quad A_{ff} = (ab)^2, \qquad A_{\Delta\Delta} = (\Delta\Delta_1)^2, \qquad A_{f\Delta} = (a\Delta)^2,
\]
\[
\qquad F = (ap)^2, \qquad M = -(\Theta p)^2.
\]
Die Ueberschiebungen dieser Formen geben zu zahlreichen Relationen Veranlassung, welche bei den verschiedenen Aufgaben, in denen gerade dieses simultane System eine Rolle spielt, von Wichtigkeit werden. Wir gelangen zu diesen Relationen einmal, indem wir frühere allgemeine Theorien über simultane Systeme oder Functionaldeterminanten benutzen, dann durch directes Ueberschieben, insbesondere der quadratischen über die linearen Formen, und der linearen unter sich, endlich aber auch, indem wir versuchen, ein und dieselbe Form auf mehrere Arten darzustellen. Jedem dieser drei Wege entsprechend, wollen wir nun eine Reihe von Relationen ermitteln.

301. \textit{Erste Reihe von Relationen.} Das simultane System von $f$ und $\Delta$, welches aus den sechs Formen
\[
f, \quad \Delta, \quad \Theta, \quad A_{ff}, \quad A_{\Delta\Delta}, \quad A_{f\Delta}
\]
besteht, liefert sofort eine Reihe von Relationen, die wir schon in der Theorie simultaner quadratischer Formen aufgestellt haben. Wir erhielten damals (vgl. Nr. 127):
\weq{I}{\left.\begin{aligned} 2A_{\Theta\Theta} &= A_{ff}A_{\Delta\Delta} - A_{f\Delta}^2 \\ A_{f\Theta} &= 0, \quad A_{\Delta\Theta} = 0 \end{aligned}\right\}}
\weq{II}{\left.\begin{aligned} (1) &\quad 2(\Theta, \; f) = A_{ff}\Delta - A_{f\Delta}f \\ (2) &\quad 2(\Delta, \; \Theta) = A_{\Delta\Delta}f - A_{f\Delta}\Delta \\ (3) &\quad -2\Theta^2 = A_{ff}\Delta^2 - 2A_{f\Delta}f\Delta + A_{\Delta\Delta}f^2 \end{aligned}\right\}}
Die letzte dieser Relationen liefert sofort die in Nr. 299 erwähnte Invarianten-Beziehung; denn ersetzen wir in ihr $x$ durch $p$ und setzen einstweilen
\[
(\Delta p)^2 = L,
\]
%CONT>>


% ===== pdf idx 339, printed 325.0 =====
% !PRINTED 325
\weberpage{325}
so kommt:

\weq{III}{-\,2\,M^2 = A_{ff}L^2 - 2\,A_{f\Delta}F\cdot L + A_{\Delta\Delta}F^2\,.}
Wie wir hier das Quadrat der beiden schiefen Formen $\Theta$ und $M$ durch
gerade Formen dargestellt haben, so können wir auch die Quadrate der drei
anderen schiefen Formen, $q$, $r$, $s$ des Systemes durch gerade Formen
darstellen. Man findet nämlich in einfacher Weise

\weq{IV}{\left.\begin{aligned}
(1)&\quad q^2 = f\cdot F - \tfrac{1}{2}\,p^2\,A_{ff}\\[6pt]
(2)&\quad r^2 = \Delta\cdot L - \tfrac{1}{2}\,p^2\,A_{\Delta\Delta}\\[6pt]
(3)&\quad -\,s^2 = \Theta\cdot M + \tfrac{1}{2}\,p^2\,A_{\Theta\Theta}
\end{aligned}\right\}\,,}
indem man entweder die früher aufgestellten Sätze über
Functionaldeterminanten zu Grunde legt, oder noch kürzer, indem man
Identitäten wie $p_x\,(ab) = a_x\,(bp) - b_x\,(ap)$ und die beiden analogen
quadrirt.

302. \textit{Zweite Reihe von Relationen.} Wir können ferner
Formenbeziehungen aufstellen durch directe Ueberschiebung der Formen $p$,
$q$, $r$, $s$ unter sich oder über die Formen $f$, $\Delta$, $\Theta$. Sie
sind durchwegs leicht zu berechnen, und wir können uns daher auf solche
Beispiele beschränken, die wir später bei der typischen Darstellung von $f$
und $\varphi$ verwenden werden. Das sind zunächst die Ueberschiebungen von
$f$ über die linearen Covarianten. Man findet:

\weq{V}{\left.\begin{aligned}
(f,\ p) &= \phantom{-}q \quad \text{(vgl. Nr. 300)}\\[6pt]
(f,\ q) &= -\tfrac{1}{2}\,A_{ff}\,p\\[6pt]
(f,\ r) &= -\,s + \tfrac{1}{2}\,A_{f\Delta}\,p\,^{*)}\\[6pt]
(f,\ s) &= +\tfrac{1}{2}\left\{A_{f\Delta}\,q + A_{ff}\,r\right\}
\end{aligned}\right\}.}
Zur Verification der drei letzten Relationen dieser Gruppe mögen die
Entwicklungen dienen:

\[
\begin{aligned}
(f,\ q) &= (f,\ (f,\ p)) = (ab)(bp)a_x = \tfrac{1}{2}\,(ab)\left\{(bp)a_x -
(ap)b_x\right\}\\
&= -\tfrac{1}{2}\,(ab)^2\cdot p\\
(f,\ r) &= (f,\ (p,\ \Delta)) = (p\Delta)(a\Delta)a_x\,;\quad
(a\Delta)a_x\Delta_y = \Theta_x\Theta_y + \tfrac{1}{2}\,A_{f\Delta}(xy)\\
&\quad\text{[vgl. Nr. 125 (2)]}\\
(f,\ s) &= ((f,\ (\Theta,\ p))) = (\Theta p)(a\Theta)a_x = -((\Theta,f),p) =
-\tfrac{1}{2}(A_{ff}\Delta - A_{f\Delta}f,\ p)\,.
\end{aligned}
\]

\par\noindent\rule{0.3\linewidth}{0.4pt}\par
{\footnotesize $^{*)}$ In Clebsch, \glqq Bin\"are Formen\grqq{} ist $(f,\
r) = +\,s$.}


% ===== pdf idx 340, printed 326.0 =====
% !PRINTED 326
\weberpage{326}
Wir fügen hinzu:

\weq{VI}{\left.\begin{aligned}
(qp) &= ((f,\ p),\ p) = (ap)^2 = F\\
(rp) &= -\,((\Delta,\ p),\ p) = -\,(\Delta p)^2 = -\,L\\
(sp) &= ((\Theta,\ p),\ p) = (\Theta p)^2 = -\,M\\ % checked against the GDZ scan (27 Oct 2026): the round dot above the right end of the "=" before -M is an ink speck; plain "=" (the old note sat inside \text and was printed)
(qr) &= ((f,\ p),\ (p,\ \Delta)) = (-(f,\ \Delta),\ p^2)^2 = -\,(\Theta p)^2 = +\,M\\
(qs) &= ((f,\ p),\ (\Theta,\ p)) = ((f,\ \Theta),\ p^2)^2 = -\tfrac{1}{2}\left\{A_{ff}L - A_{f\Delta}F\right\}\\
(rs) &= ((\Delta,\ p),\ (\Theta,\ p)) = ((\Delta,\ \Theta),\ p^2)^2 = -\tfrac{1}{2}\left\{A_{f\Delta}L - A_{\Delta\Delta}F\right\}
\end{aligned}\right\}.}
Mit Hilfe dieser beiden Gruppen (V) und (VI) von Relationen lassen sich nun
auch sofort die Werthe jener Invarianten anschreiben, welche entstehen,
wenn wir die Form $f$ zweimal über irgend ein Product zweier linearer
Covarianten schieben. Man findet unmittelbar:

\weq{VII}{\left.\begin{aligned}
(f,\ pq)^2 &= (q,\ q) = 0\\
(f,\ pr)^2 &= (q,\ r) = +\,M\\
(f,\ ps)^2 &= (q,\ s) = -\tfrac{1}{2}\left\{A_{ff}L - A_{f\Delta}F\right\}\\
(f,\ qr)^2 &= -\tfrac{1}{2}\,A_{ff}(p,\ r) = -\tfrac{1}{2}\,A_{ff}L\\
(f,\ qs)^2 &= -\tfrac{1}{2}\,A_{ff}(p,\ s) = -\tfrac{1}{2}\,A_{ff}M\\
(f,\ rs)^2 &= \left(-\,s + \tfrac{1}{2}\,A_{f\Delta}p,\ s\right) = \tfrac{1}{2}\,A_{f\Delta}M
\end{aligned}\right\}.}
303. \textit{Dritte Reihe von Relationen.} Die dritte Reihe von Relationen
hat ihren Ursprung in der Fundamentalrelation dieses simultanen Systemes

\weq{VIII}{\Theta = (f,\ \Delta) = (\varphi,\ p)\,.}
Wir beweisen dieselbe, indem wir die Form

\[
\Theta = (a\Delta)\,a_x\,\Delta_x = (\alpha\beta)^2(a\alpha)\,a_x\,\beta_x
\]
mit Hilfe des Identitätssatzes umformen. Man hat nämlich

\[
\begin{aligned}
\Theta &= (\alpha\beta)^2(a\alpha)\,a_x\,\beta_x = (a\alpha)(\alpha\beta)\,\beta_x\left\{\alpha_x(a\beta) - \beta_x(a\alpha)\right\}\\
&= (a\alpha)(a\beta)(\alpha\beta)\,\alpha_x\,\beta_x - (a\alpha)^2(\alpha\beta)\,\beta_x^2\\
&= -\,(a\alpha)^2(\alpha\beta)\,\beta_x^2\,,
\end{aligned}
\]
da das erste Glied rechts identisch verschwindet. Führen wir aber in
$(a\alpha)^2(\alpha\beta)\,\beta_x$ das Symbol $p_x = (a\alpha)^2\alpha_x$
ein, so kommt:

\[
\Theta = -\,(p\beta)\,\beta_x^2 = (\varphi,\ p)\,.
\]
Aus dieser fundamentalen Relation (VIII) gehen insbesondere folgende
wichtige Beziehungen hervor:
%CONT>>


% ===== pdf idx 341, printed 327.0 =====
% !PRINTED 327
\weberpage{327}
%<<CONT
\weq{IX}{\left.\begin{aligned}
s &= (\Theta,\ p) = ((\varphi,\ p),\ p) = (\varphi,\ p)^2 = (\alpha p)^2\alpha_x\\
M &= -\,(s,\ p) = -\,((\varphi,\ p)^2,\ p) = -\,(\alpha p)^3\\
-\,r &= (\Delta p)\,\Delta_x = (\alpha\beta)^2(\beta p)\,\alpha_x^2 = (\alpha_x^3,\ (\beta p)\,\beta_x)^2 = (\varphi,\ \Theta)^2\\
L &= (\Delta p)^2 = (\alpha\beta)^2(\alpha p)(\beta p) = ((\alpha p)\,\alpha_x^2,\ (\beta p)\,\beta_x^2) = A_{\Theta\Theta}\\
&= \tfrac{1}{2}\left\{A_{ff}A_{\Delta\Delta} - A_{f\Delta}^2\right\}\\
(\Theta,f) &= p^2 + (\varphi,\ q)
\end{aligned}\right\}.}
Was die letzte dieser Relationen betrifft, so ergiebt sich ihre
Richtigkeit aus folgender Entwicklung. Es ist

\[
\begin{aligned}
(\Theta,\ f) &= ((\varphi,\ p),\ f) = ((\alpha p)\,\alpha_x,\ a_x^2) = (\alpha p)(\alpha a)\,\alpha_x\,a_x\\
&= (\alpha a)\,\alpha_x\left\{\alpha_x(ap) + p_x(\alpha a)\right\}\\
&= (\alpha q)\,\alpha_x^2 + p_x^2 = (\varphi,\ q) + p^2\,.
\end{aligned}
\]
Wir wollen damit die Untersuchungen über die Covariantenbeziehungen des
simultanen Systemes von $f$ und $\varphi$ beschliessen und nun zu einigen
Anwendungen übergehen.

304. \textit{Typische Darstellung.} Die beiden Formen $f$ und $\varphi$
lassen sich mit Hilfe zweier ihrer vier simultanen linearen Invarianten
typisch darstellen. Gehen wir zu dem Zwecke wiederum von Identitäten wie

\[
\begin{aligned}
(pq)^2\,a_x^2 &= \left\{(aq)\,p_x - (ap)\,q_x\right\}^2\\
(pq)^3\,\alpha_x^3 &= \left\{(\alpha q)\,p_x - (\alpha p)\,q_x\right\}^3
\end{aligned}
\]
aus, so ergeben sich für die erste derselben die Werthe der typischen
Coefficienten unmittelbar aus den Relationen in Nr. 302. Man hat
beispielsweise:

\[
\begin{aligned}
F^2\cdot f &= (aq)^2\cdot p_x^2 - 2\,(ap)(aq)\,p_x\,q_x + (ap)^2\,q_x^2\\
&= (f,\ q)^2\cdot p_x^2 - 2\,(f,\ pq)\,p_x\,q_x + (f,\ p)^2\cdot q_x^2\\
&= \tfrac{1}{2}\,A_{ff}F\cdot p_x^2 + F\cdot q_x^2
\end{aligned}
\]
oder

\[
F\cdot f = \tfrac{1}{2}\,A_{ff}\,p_x + q_x^2\,,
\]
ein Resultat, das wir auch schon an anderer Stelle (Nr. 301, IV) ermittelt
haben.

Die typischen Coefficienten für die zweite Identität

\[
-\,F^3\cdot\varphi = (\alpha q)^3\,p_x^3 - 3\,(\alpha q)^2(\alpha p)\,p_x^2\,q_x + 3\,(\alpha q)(\alpha p)^2\,p_x\,q_x^2 - (\alpha p)^3\,q_x^3
\]
berechnet man mit Hilfe der in Nr. 303 entwickelten Relationen:

\[
\begin{aligned}
\Theta &= (\varphi,\ p)\ = (\alpha p)\,\alpha_x^2\\
s &= (\varphi,\ p)^2 = (\alpha p)^2\,\alpha_x\,.
\end{aligned}
\]


% ===== pdf idx 342, printed 328.0 =====
% !PRINTED 328
\weberpage{328}
Aus diesen ergiebt sich sofort:

\[
\begin{aligned}
(\alpha p)^3 &= (sp) = -\,M\\
(\alpha p)^2(\alpha q) &= (sq) = \tfrac{1}{2}\left\{A_{ff}L - A_{f\Delta}F\right\}\\
(\alpha p)(\alpha q)^2 &= (\Theta,\ q^2)^2 = \left(\Theta,\ F\cdot f - \tfrac{1}{2}\,A_{ff}p^2\right)^2 = -\tfrac{1}{2}\,A_{ff}(\Theta p)^2 = \tfrac{M}{2}\cdot A_{ff}\\
(\alpha q)^3 &= \left((\alpha q)\,\alpha_x^2,\ q^2\right)^2 = \left((\alpha q)\,\alpha_x^2,\ Ff - \tfrac{1}{2}\,A_{ff}p^2\right)^2\\
&= F\cdot(\alpha q)(\alpha a)^2 - \tfrac{1}{2}\,A_{ff}\cdot(\alpha q)(\alpha p)^2 = F\cdot(pq) - \tfrac{1}{2}\,A_{ff}(sq)\,.
\end{aligned}
\]
In ebenso einfacher Weise würden sich die typischen Coefficienten
berechnen lassen, wenn wir $f$ und $\varphi$ durch irgend zwei andere
lineare Covarianten hätten darstellen wollen.

305. \textit{Ausnahmefälle für typische Darstellung.} Die typische
Darstellung wird unmöglich, sobald die beiden Invarianten

\[
\begin{aligned}
F &= (ap)^2\\
L &= (\Delta p)^2
\end{aligned}
\]
gleichzeitig identisch verschwinden; denn in diesem Falle verschwinden
auch alle Invarianten

\[
(pq),\ (pr),\ (ps),\ (qr),\ (qs),\ (rs)\,,
\]
wie ein Blick auf die Relationen (VI) Nr. 302 lehrt, und demnach sind die
vier linearen Covarianten nicht mehr von einander unabhängig.

Das Verschwinden von $F$ und $L$ sagt aber nichts anderes aus, als dass
$f$ und $\Delta$ den gemeinsamen Factor $p$ besitzen. In diesem Falle muss
aber $\varphi$ den linearen Factor $p$ quadratisch enthalten. Denn
multipliciren wir

\weq{1}{(ap)\,\Delta_x - (\Delta p)\,a_x = (a\Delta)\,p_x}
mit

\weq{2}{a_y\,\Delta_x + \Delta_y\,a_x = 2\,(a_x\,\Delta_x)_y\,,}
so kommt:

\[
\bigl((ap)\,\Delta_x - (Ap)\,a_x\bigr)\,\bigl(a_y\,\Delta_x + \Delta_y\,a_x\bigr) = 2\,\bigl((a\Delta)\,a_x\,\Delta_x\bigr)_y\,p_x\,. % [sic] \textquotedbl{}(Ap)\textquotedbl{} printed with roman A, not Delta, as elsewhere in this line
\]
Ersetzen wir aber hierin $y$ durch $p$, so erhalten wir

\[
(ap)^2\,\Delta_x^2 - (\Delta p)^2\,a_x^2 = 2\,p_x\,(\Theta,\ p)\,.
\]
Die linke Seite ist null wegen $F = L = 0$; also ist auch

\[
(\Theta,\ p) = (\alpha p)^2\,\alpha_x = 0\,,
\]
d.\,h. $\varphi$ besitzt den Factor $p$ im zweiten Grade.

Ebenso findet man umgekehrt, dass, wenn $\varphi$ einen linearen Factor
von $f$ quadratisch besitzt, $F$ und $L$ verschwinden müssen und sonach
$p$, $q$, $r$, $s$ einander proportional sind.


% ===== pdf idx 343, printed 329.0 =====
% !PRINTED 329
\weberpage{329}

$F$ und $L$ verschwinden auch, wenn $\Delta$ proportional $f$ ist, also
wenn $\Theta = 0$; dann kann aber überhaupt nicht mehr von einem simultanen
System von $f$ und $\varphi$ die Rede sein.

306. \textit{Resultante von f und $\varphi$.} Um die Resultante von $f$ und
$\varphi$ aufzustellen, benutzen wir die nämlichen Principien, wie wir sie
beim Aufsuchen der Discriminanten der Form fünfter und sechster Ordnung
verwendeten. Es möge $\lambda$ der gemeinsame Factor von $f$ und $\varphi$
sein. Diesen Factor besitzt alsdann insbesondere auch die erste
Ueberschiebung von $\varphi$ über $\lambda$. Sei $\mu$ der zweite lineare
Factor von $f$, $\nu$ der zweite von $(\varphi, \lambda)$, so können wir
unsere Bedingungen formuliren:
\[
a_x^2 = \lambda\mu, \quad \alpha_x^2(\alpha\lambda) = \lambda\nu.
\]
Versuchen wir auf Grund dieser Annahmen die Invarianten des Systemes zu
berechnen, so wird sich zwischen denselben eine Relation ergeben, die
gemäss der Voraussetzung und nach dem Grade in den Coefficienten nur die
Resultante von $f$ und $\varphi$ sein kann.

Wir berechnen zunächst die Invariante $F = (ap)^2$. Die Form $p$ war
dargestellt durch:
\[
p = (\alpha a)^2 \alpha_x, \quad \text{oder weil } a_x^2 = \lambda\mu,
\]
\[
p = (\alpha\lambda)(\alpha\mu)\alpha_x = ((\alpha\lambda)a_x^2, \mu) = (\lambda\nu, \mu).
\]
Nun ist:
\[
(\lambda\nu, \mu) = \frac{1}{2}\nu(\lambda\mu) + \frac{1}{2}\lambda(\nu\mu).
\]
Da aber
\[
\nu(\lambda\mu) = (\nu\mu)\lambda + (\lambda\nu)\mu,
\]
so kommt:
\[
(\lambda\nu, \mu) = \lambda(\nu\mu) + \frac{1}{2}\mu(\lambda\nu) = p.
\]
Bilden wir somit die zweite Ueberschiebung $(ap)^2 = F$, so erhalten wir
\[
F = \left(a_x^2, \ \lambda^2(\nu\mu)^2 + \lambda\cdot\mu(\nu\mu)(\lambda\nu) + \frac{1}{4}\mu^2(\lambda\nu^2)\right)^2,
\]
oder weil $a_x^2 = \lambda\cdot\mu$, und demnach das erste und dritte
Ueberschiebungsglied gleich null, so wird:
\weq{1}{F = (a_x^2, \ b_x^2(\nu\mu)(\lambda\nu))^2 = (\nu\mu)(\lambda\nu)\cdot A_{ff}.}
Berechnen wir sodann die simultane Invariante $A_{f\Delta}$, so findet man:
\[
A_{f\Delta} = (\Delta\lambda)(\Delta\mu) = (\alpha\beta)^2(\beta\lambda)(\alpha\mu).
\]
Der Ausdruck rechts ist nichts anderes als die zweite Ueberschiebung
\[
((\beta\lambda)\beta_x^2, \ (\alpha\mu)\alpha_x^2)^2 = (\lambda\nu, \ (\alpha\mu)\alpha_x^2)^2 = (\alpha\lambda)(\alpha\nu)(\alpha\mu).
\]


% ===== pdf idx 344, printed 330.0 =====
% !PRINTED 330
\weberpage{330}

Nun ist aber:
\[
(\alpha\lambda)(\alpha\nu)(\alpha\mu) = ((\alpha\lambda)a_x^2, \ \mu\nu)^2 = (\lambda\nu, \ \mu\nu) = \frac{1}{2}(\nu\mu)(\lambda\nu),
\]
somit:
\weq{2}{2A_{f\Delta} = (\nu\mu)(\lambda\nu).}
Setzen wir den Werth von $(\nu\mu)(\lambda\nu)$ in die Relation (1) ein, so
kommt:
\[
F - 2A_{f\Delta}\cdot A_{ff} = 0.
\]
Diese Relation besteht also unter der Voraussetzung, dass $f$ und $\varphi$
einen gemeinsamen Factor haben. Sie ist überdies vom Grade $2 + 3 = 5$ in
den Coefficienten, und ihre linke Seite darf somit als Resultante von $f$
und $\varphi$ betrachtet werden.

\begin{center}
\S~32. \\
\textbf{Das simultane System zweier cubischer Formen.}
\end{center}

307. \textit{Erste Formulirung der Aufgabe.} Das vollständige System einer
cubischen Form besteht aus den vier Formen:
\[
a_x^3 = f, \quad (f,f)^2 = \Delta, \quad (f,\Delta) = Q, \quad (\Delta,\Delta)^2 = A_{\Delta\Delta}.
\]
Es sei $\varphi$ eine zweite cubische Form; die Formen ihres vollen
Systemes wollen wir bezeichnen mit
% checked against the GDZ scan (27 Oct 2026): the sign is an inverted slanted Delta (thick left stroke, thin diagonal, closed top bar), the \nabla of pp. 337-340; transcribed \nabla on pp. 330-336
\[
\alpha_x^3 = \varphi, \quad (\varphi,\varphi)^2 = \nabla, \quad (\varphi,\nabla) = K, \quad (\nabla,\nabla)^2 = A_{\nabla\nabla}.
\]
Wir stellen uns die Aufgabe, das simultane System der beiden Formen $f$ und
$\varphi$ zu ermitteln. Nach den allgemeinen Sätzen, § 38, erhalten wir
dasselbe, wenn wir die beiden Producte
\[
U = f^{\lambda_1}\cdot Q^{\lambda_2}\cdot \Delta^{\lambda_3}
\]
\[
V = \varphi^{\mu_1}\cdot K^{\mu_2}\cdot \nabla^{\mu_3}
\]
über einander schieben, d.\,h. also mit einander multipliciren und das
Product $U\cdot V$ auf alle möglichen Arten in bekannter Weise falten.
Enthält eine solche Ueberschiebung
\[
F = (U,V)^{\varrho} = (f^{\lambda_1}\cdot Q^{\lambda_2}\cdot \Delta^{\lambda_3}, \ \varphi^{\mu_1}\cdot K^{\mu_2}\cdot \nabla^{\mu_3})^{\varrho}
\]
kein zerfallendes Glied, so liefert sie eine Form des Systemes.

308. \textit{Reduction der Aufgabe.} Die Zahl der zu faltenden Producte
$U\cdot V$ lässt sich von vornherein reduciren, insofern sich bis auf
wenige Ausnahmen leicht erkennen lässt, welche derselben auf
Ueberschiebungen mit zerfallenden Gliedern führen.

Erstens: Producte, welche höhere als die zweite Potenz von $\Delta$ resp.
$\nabla$ enthalten, bleiben unberücksichtigt, wegen der bekannten Relation:
\[
\Delta^3 + 2Q^2 + A_{\Delta\Delta}f^2 = 0.
\]
Zweitens: Enthält $U$ eine Form $\Delta$ und gleichzeitig $V$ eine Form
$\nabla$, so sind höhere Faltungen als die zweite ausgeschlossen, da stets
Glieder mit dem Factor $A_{\Delta\nabla} = (\Delta,\nabla)^2$ gebildet
werden können.


% ===== pdf idx 345, printed 331.0 =====
% !PRINTED 331
\weberpage{331}

Drittens: Bei Ueberschiebungen, welche den Grad 3 übersteigen, kann $U$ nur
Einen Factor $f$ oder $Q$, und $V$ nur Einen Factor $\varphi$ oder $K$
enthalten, da diese Formen von gleichem Grade sind und somit stets Glieder
auftreten würden, welche die dritten Ueberschiebungen $(f,\varphi)^3$,
$(f,K)^3$, $(\varphi,Q)^3$ zum Factor haben. Bei niedrigeren
Ueberschiebungen ist aber ohnehin der Fall ausgeschlossen, dass $U$ oder
$V$ aus mehr als aus einem Factor besteht, wenn dieser dritten Grades ist.

Viertens: Ueberschiebungen einer Potenz von $f$, oder $Q$, über eine Potenz
von $\nabla$, bleiben unberücksichtigt, da schon $(f^2,\nabla^2)^4$, oder
$(Q^2,\nabla^2)^4$ Glieder enthalten, welche Producte früherer
Ueberschiebungen sind, nämlich: $(f,\nabla)^2\cdot(f,\nabla)^2$ oder
$(Q,\nabla)^2\cdot(Q,\nabla)^2$.

Von allen Producten $U\cdot V$ sind also nur jene zu falten, in welchen\\
1) sowohl $U$ als $V$ nur \gesperrt{eine} Form enthält,\\
2) $U$ (oder $V$) eine Form, $V$ (oder $U$) eine zweite Potenz von
$\nabla$ (resp. $\Delta$).

309. \textit{Aufstellung der einzelnen Formen.} Wir bilden nun die
einzelnen Ueberschiebungen:
\[
1) \quad U = f, \quad V = \varphi.
\]
Ihre Ueberschiebung liefert drei Formen:
\[
(f,\varphi) = \vartheta, \quad (f,\varphi)^2 = \Theta, \quad (f,\varphi)^3 = J.
\]
\[
2) \quad U = f, \quad V = \nabla, \quad \text{oder} \quad U = \Delta, \quad V = \varphi.
\]
Wir erhalten sechs Formen:
\[
\begin{aligned}
&(f,\nabla), \ (f,\nabla)^2, \ (f,\nabla^2)^3\\
&(\varphi,\Delta), \ (\varphi,\Delta)^2, \ (\varphi,\nabla^2)^3.
\end{aligned}
\]
\[
3) \quad U = \Delta, \quad V = \nabla.
\]
Es entstehen die zwei Formen:
\[
(\Delta,\nabla) \quad \text{und} \quad (\Delta,\nabla)^2.
\]
\[
4) \quad U = Q \quad \text{und} \quad V
\]
sei irgend eine Form des zweiten Systems. Als erste Ueberschiebungen
treten die Formen auf:
\[
(Q,\varphi), \ (Q,\nabla); \ (Q,K),
\]
und entsprechend
\[
(K,f), \ (Q,\Delta);
\]
sie sind als Functionaldeterminanten von Functionaldeterminanten
reducibel.

Als zweite Ueberschiebungen erhält man:
\[
\begin{aligned}
&(Q,\varphi)^2, \ (Q,\nabla)^2, \ (Q,K)^2\\
&(K,f)^2, \ (Q,\Delta)^2.
\end{aligned}
\]


% ===== pdf idx 346, printed 332.0 =====
% !PRINTED 332
\weberpage{332}

Von diesen ist $(Q,K)^2$ reducibel; denn falten wir das Product
\[
a_x^2\Delta_x(a\Delta)\cdot \alpha_x^2\nabla_x(\alpha\nabla)
\]
zweimal, und bilden insbesondere das Glied
\[
G = a_x\alpha_x(a\nabla)(\alpha\nabla)(a\Delta)(\alpha\Delta),
\]
% checked against the GDZ scan (27 Oct 2026): the dot above the a of the first (a\nabla) is an ink speck (the same product in the next display has a plain a); no accent
so lässt sich auf dasselbe der Productsatz unmittelbar anwenden und man
findet
\[
G = \alpha_x\alpha_x(a\nabla)(\alpha\nabla)(a\Delta)(\alpha\Delta) % [sic]
\]
\[
= \frac{1}{2}a_x\alpha_x\{(\nabla\alpha)^2(\Delta a)^2 + (\nabla a)^2(\Delta\alpha)^2 - (\nabla\Delta)^2(a\alpha)^2\}.
\]
Das erste Glied rechts ist wegen $(\Delta a)^2$ oder $(\nabla\alpha)^2$
null. [Vergl. Nr. 145, (I).] Das zweite ist $(f,\nabla)^2\cdot(\varphi,\Delta)^2$,
das dritte $(\nabla,\Delta)^2\cdot(f,\varphi)^2$. Es ist somit das Glied
$G$ reducibel, folglich $(Q,K)^2$ selbst.

Als dritte Ueberschiebungen erhält man endlich:
\[
\begin{aligned}
&(Q,\varphi)^3, \ (Q,f)^3, \ (Q,K)^3\\
&(Q,\nabla^2)^3, \ (K,\nabla^2)^3.
\end{aligned}
\]
Von diesen Ueberschiebungen sind die beiden letzten reducibel. Es ist nur
nöthig, dies von $(K,\Delta^2)^3$ zu zeigen; für die Form $(Q,\nabla^2)^3$
folgt die Reducibilität aus der entsprechenden Bedeutung der Formen $Q$
und $\nabla$. Die dritte Polare von $K$ ist:
\[
K_{y^3} = (\alpha\nabla)\alpha_y^2\nabla_y,
\]
ferner die dritte Polare von $\Delta_x^4$
\[
\Delta_{y^3}^2 = \Delta_{y^2}\Delta_{1y}\Delta_{1x}.
\]
Indem wir das Product
\[
(\alpha\nabla)\alpha_y^2\nabla_y\cdot \Delta_{y^2}\Delta_{1y}\Delta_{1x}
\]
auf alle möglichen Arten dreimal falten, erhalten wir die Ueberschiebung
\[
(K,\Delta^2)^3 = \frac{1}{3}(\alpha\nabla)(\alpha\Delta)^2(\nabla\Delta_1)\Delta_{1x} + \frac{2}{3}(\alpha\nabla)(\alpha\Delta_1)(\nabla\Delta)(\alpha\Delta).
\]
Die Differenz dieser beiden Glieder ist aber null, da sie den Factor
$(\alpha\nabla)^2$ enthält (vergl. Nr. 145), folglich sind beide einander
gleich, und wir erhalten:
\[
(K,\Delta^2)^3 = (\alpha\nabla)(\alpha\Delta)^2(\nabla\Delta_1)\Delta_{1x}.
\]
Ersetzen wir hier $\Delta_1$ durch $(ab)^2a_xb_x$, so kommt:
\[
(K,\Delta^2)^3 = (ab)^2(\alpha\Delta)^2(\alpha\nabla)(\nabla a)b_x
\]
\[
= (\alpha\nabla)(\nabla a)b_x\{(b\Delta)(a\alpha) - (a\Delta)(b\alpha)\}^2.
\]
Entwickelt man rechts nach dem binomischen Lehrsatz, so ist das erste und
letzte Glied null wegen des Factors $(b\Delta)^2$. Es reducirt sich also:
%CONT>>


% ===== pdf idx 347, printed 333.0 =====
% !PRINTED 333
\weberpage{333}

\[
(K, \Delta^2)^3 = -2(\alpha\nabla)(\nabla a)(a\alpha)(b\alpha)(a\Delta)(b\Delta)b_x
\]
\[
= -2(\alpha\nabla)(a\Delta)(b\Delta)(a\alpha)b_x\{(\nabla b)(a\alpha)-(ab)(\nabla\alpha)\}
\]
\[
= -2(a\alpha)^2(\alpha\nabla)(\nabla b)(a\Delta)(b\Delta)b_x.
\]
% checked against the GDZ scan (27 Oct 2026): the sign paired with \Delta on pp. 330-336 is an
% inverted slanted Delta (thick left stroke, thin diagonal, closed top bar), the same sign as
% the \nabla of pp. 337-340; the italic V of this fount (p. 234) has open top serifs and
% no closed top. Transcribed \nabla on pp. 330-336.

Da rechts der Klammerfactor $(a\alpha)^2$ auftritt, so suchen wir diesen
Ausdruck mittelst der Reihenentwicklung
\[
(a\alpha)^2 a_y\alpha_x = \Theta_x\Theta_y + \frac{1}{2}J(yx)
\]
umzuformen. Ersetzt man nämlich in ihr $y$ durch $\Delta$, $x$ durch $\nabla$,
so kommt:
\[
(a\alpha)^2(a\Delta)(\alpha\nabla) = (\Theta\nabla)(\Theta\Delta) + \frac{1}{2}J(\Delta\nabla).
\]

Tragen wir dies in $(K, \Delta^2)^3$ ein, so erhalten wir:
\[
(K, \Delta^2)^3 = -2b_x(b\Delta)(\nabla b)\{(\Theta\nabla)(\Theta\Delta)+\frac{1}{2}J(\Delta\nabla)\}
\]
\[
= +2b_x(b\Delta)(\Theta\Delta)(b\nabla)(\Theta\nabla) + J(\Delta\nabla)(b\Delta)(b\nabla)b_x.
\]
Für das erste Glied liefert der Productsatz:
\[
(b\Delta)(\Theta\Delta)(b\nabla)(\Theta\nabla) = \frac{1}{2}\{(b\Delta)^2(\Theta\nabla)^2+(\Theta\Delta)^2(b\nabla)^2-(\Delta\nabla)^2(b\Theta)^2\}.
\]

Das erste Glied in dieser Summe ist null wegen $(b\Delta)^2$. Man erhält
also:
\[
(K, \Delta^2)^3 = (\Theta, \Delta)^2\cdot(f, \nabla)^2 - (\Delta, \nabla)^2\cdot(f, (f, \varphi))^2 + J\cdot(f, (\Delta, \nabla))^2.
\]

Die beiden Covarianten $(K\Delta^2)^3$ und $(Q\nabla^2)^3$ sind somit
reducibel und das ganze Formensystem der zwei cubischen Formen $f$ und
$\varphi$ besteht also aus 26 Formen.

Eine biquadratische:
\[
(f, \varphi) = (a\alpha)a_x^2\alpha_x^2 = \vartheta.
\]

Sechs cubische:
\[
f,\ \varphi,\ Q,\ K,\ (f, \nabla),\ (\varphi, \Delta).
\]

Sechs quadratische:
\[
\Delta,\ \nabla,\ (f, \varphi)^2 = \Theta,\ (\Delta, \nabla),\ (Q, \varphi)^2,\ (K, f)^2.
\]

Sechs lineare:
\[
(f, \nabla)^2,\ (\varphi, \Delta)^2,\ (f, \nabla^2)^3,\ (\varphi, \Delta^2)^3,\ (Q, \nabla)^2,\ (K, \Delta)^2.
\]

Sieben Invarianten:
\[
A_{\Delta\Delta},\ A_{\nabla\nabla},\ A_{\Delta\nabla},\ (f, \varphi)^3 = J,\ (Q, K)^3,\ (Q, \varphi)^3,\ (Q, f)^3.
\]
Von den linearen Formen erhalten zwei eine besondere Bezeichnung.
Wir setzen
\[
(f, \nabla)^2 = (a\nabla)^2a_x = \pi_x = \pi
\]
\[
(\varphi, \Delta)^2 = (\alpha\Delta)^2\alpha_x = p_x = p.
\]


% ===== pdf idx 348, printed 334.0 =====
% !PRINTED 334
\weberpage{334}

310. \textit{Exposition der folgenden Aufgaben.} Nachdem wir im
Vorhergehenden die Formen ermittelt haben, aus denen das vollständige
System von $f$ und $\varphi$ bestehen muss, wird unser nächstes Ziel nach
den Beziehungen und Relationen gerichtet sein, welche, sei es zwischen
den Formen des Systemes selbst, sei es zwischen deren Ueberschiebungen
bestehen. Die allgemeinen Wege, welche dazu führen, sowie die hiebei
sich einstellende Nothwendigkeit der Einführung neuer gleich
berechtigter Formen an Stelle alter Formen, haben wir wiederholt an
den geeigneten Stellen klar gelegt. Wir greifen hier indes der
systematischen Methode vor, indem wir uns zunächst, in Folge der
Wichtigkeit der sich hiebei einstellenden Resultate, mit den linearen
Covarianten $p$ und $\pi$ beschäftigen, deren erste Ueberschiebung über
$f$ und $\varphi$ sowohl, als deren Quadrate wir vor allem berechnen
wollen. Hiebei ist es, wie überhaupt bei allen kommenden
Untersuchungen, stets hinreichend, die Berechnungen für eine je zweier
gleich berechtigter Formen $(f, p)$ und $(\varphi, \pi)$, oder $(f, \pi)$
und $(\varphi, p)$, oder $p^2$ und $\pi^2$ durchzuführen, da die
Resultate für die andere der beiden Formen einfach durch Vertauschung
von $a$ mit $\alpha$ gewonnen werden.

311. \textit{Die ersten Ueberschiebungen von $f$ und $\varphi$ über $p$
und $\pi$.} Die erste Ueberschiebung von $f$ über $p$ ist dargestellt
durch:
\[
(f, p) = a_x^2(a\alpha)(\alpha\Delta)^2.
\]
Subtrahiren wir auf der rechten Seite die Grösse
\[
0 = (a\alpha)(a\Delta)^2\alpha_x^2,
\]
so erhalten wir
\[
(f, p) = (a\alpha)\{a_x^2(\alpha\Delta)^2 - \alpha_x^2(a\Delta)^2\}
\]
\[
= (a\alpha)\{a_x(\alpha\Delta) + \alpha_x(a\Delta)\}\{a_x(\alpha\Delta) - \alpha_x(a\Delta)\}.
\]
Der letzte Factor rechts ist aber gleich $-(a\alpha)\Delta_x$ gemäss dem
Identitätssatze. Folglich ist:
\[
(f, p) = -(a\alpha)^2\Delta_x\{a_x(\alpha\Delta) + \alpha_x(a\Delta)\},
\]
oder:
\weq{1}{(f, p) = -2\big((a\alpha)\big)^2a_x\alpha_x, \Delta\big) = -2(\Theta, \Delta).}
% checked against the GDZ scan (27 Oct 2026): the print has -2((a\alpha))^2 a_x\alpha_x, \Delta), with a
% bold "(" before (a\alpha), a bold ")" carrying the exponent 2 directly after (a\alpha), and a
% bold ")" after \Delta; reproduced as printed (meant: -2((a\alpha)^2 a_x\alpha_x, \Delta)).
Also ebenso:
\weq{2}{(\varphi, \pi) = -2(\Theta, \nabla).}
Es ist ferner:
\[
(f, \pi) = a_x^2(ab)(b\nabla)^2.
\]
Vertauscht man wegen des Factors $(ab)$ das Symbol $a$ mit $b$ und
nimmt die halbe Summe, so kommt:
\[
(f, \pi) = \frac{1}{2}(ab)\{a_x^2(b\nabla)^2 - b_x^2(a\nabla)^2\}
\]
\[
= -\frac{1}{2}(ab)^2\nabla_x\{a_x(b\nabla) + b_x(a\nabla)\},
\]
%CONT>>


% ===== pdf idx 349, printed 335.0 =====
% !PRINTED 335
\weberpage{335}
%<<CONT
oder:
\weq{3}{(f, \pi) = -\big((ab)^2a_xb_x, \nabla\big) = -(\Delta, \nabla).}
Also auch:
\weq{4}{(\varphi, p) = -(\nabla, \Delta).}
Diese beiden Ueberschiebungen sind also bis aufs Vorzeichen einander
gleich, und man gewinnt daher aus (3) und (4) die
\gesperrt{Fundamentalrelation:}
\weq{5}{(f, \pi) + (\varphi, p) = 0.}

312. \textit{Eine Untersuchung über die Fundamentalrelation
$(f, \pi)+(\varphi, p)=0$.} Ehe wir dazu übergehen, die Quadrate von $p$
und $\pi$ zu berechnen, wollen wir uns die Frage vorlegen, ob nicht die
Formen $p$ und $\pi$ bereits durch die Relation
\weq{1}{(f, \pi) + (\varphi, p) = 0}
vollständig definirt sind. Die zu stellende Frage können wir in der
Form aussprechen: Für welche Werthe von $y$ und $z$ wird die Summe der
beiden Polaren
\[
f_y + \varphi_z
\]
identisch null? So lange $f$ und $\varphi$ Formen beliebiger Ordnung
sind, wird es nicht möglich sein, wenn nicht anderweitige Bedingungen
stattfinden, solche Werthe zu bestimmen. Nur in dem Falle, wo $f$ und
$\varphi$ quadratische Formen, giebt es unendlich viele Werthe. Bei
Formen dritter Ordnung findet man nur ein einziges Werthsystem, bei
biquadratischen überhaupt keines, wenn nicht eine bestimmte simultane
Invariante verschwindet. Dies lehrt schon eine einfache
Constantenabzählung.

Wenn wir nun hier für cubische Formen die Frage beantworten wollen, so
verlangt zunächst die Identität:
\[
f_y + \varphi_z = 0,
\]
da sie unabhängig von $x$ erfüllt sein soll, dass die Coefficienten
von $x_1^2, 2x_1x_2, x_2^2$ für sich allein verschwinden. Das liefert
drei Gleichungen ersten Grades für die Verhältnisse der vier Grössen
$y_1\ y_2\ z_1\ z_2$, nämlich:
\[
\bar a_0y_1 + \bar a_1y_2 + \bar\alpha_0z_1 + \bar\alpha_1z_2 = 0
\]
\[
\bar a_1y_1 + \bar a_2y_2 + \bar\alpha_1z_1 + \bar\alpha_2z_2 = 0
\]
\[
\bar a_2y_1 + \bar a_3y_2 + \bar\alpha_2z_1 + \bar\alpha_3z_2 = 0.
\]
Gemäss der Theorie der Auflösung linearer Gleichungen (vergl. Bd. I,
Nr. 111) sind also die Werthe von $y_1\ y_2\ z_1\ z_2$ proportional den
einzelnen Determinanten der Matrix
\[
V = \begin{vmatrix} \bar a_0 & \bar a_1 & \bar\alpha_0 & \bar\alpha_1 \\ \bar a_1 & \bar a_2 & \bar\alpha_1 & \bar\alpha_2 \\ \bar a_2 & \bar a_3 & \bar\alpha_2 & \bar\alpha_3 \end{vmatrix}.
\]


% ===== pdf idx 350, printed 336.0 =====
% !PRINTED 336
\weberpage{336}

Es ist:
\[
\varrho y_1 = \begin{vmatrix} \bar a_1 & \bar\alpha_0 & \bar\alpha_1 \\ \bar a_2 & \bar\alpha_1 & \bar\alpha_2 \\ \bar a_3 & \bar\alpha_2 & \bar\alpha_3 \end{vmatrix} = \begin{vmatrix} a_1^2a_2 & \alpha_1^3 & \beta_1^2\beta_2 \\ a_1a_2^2 & \alpha_1^2\alpha_2 & \beta_1\beta_2^2 \\ a_2^3 & \alpha_1\alpha_2^2 & \beta_2^3 \end{vmatrix} = a_2\alpha_1\beta_2(a\alpha)(a\beta)(\alpha\beta).
\]
Vertauscht man rechts $\alpha$ mit $\beta$ und nimmt die halbe Summe,
so kommt:
\[
\varrho y_1 = \frac{1}{2}a_2(\alpha\beta)^2(a\alpha)(a\beta)
\]
oder:
\[
\varrho y_1 = \frac{1}{2}a_2(a\nabla)^2 = \frac{1}{2}\bar\pi_2.
\]
Ebenso:
\[
\varrho y_2 = -\frac{1}{2}a_1(a\nabla)^2 = -\frac{1}{2}\bar\pi_1
\]
\[
\varrho z_1 = \frac{1}{2}\alpha_2(\alpha\Delta)^2 = \frac{1}{2}\bar p_2
\]
\[
\varrho z_2 = -\frac{1}{2}\alpha_1(\alpha\Delta)^2 = -\frac{1}{2}\bar p_1.
\]
% checked against the GDZ scan (27 Oct 2026): the overbars on \pi_2, \pi_1, p_2, p_1 are clearly printed;
% reproduced as printed.

Man erkennt, $p$ und $\pi$ sind die einzigen Werthe, für welche die
Identität (1) besteht, und sie sind demnach durch sie vollständig
definirt.

Bei quadratischen Formen wäre die Matrix $V$ nur zweizeilig gewesen;
zwei lineare homogene Gleichungen mit vier Unbekannten werden aber
durch unendlich viele Werthe befriedigt; bei biquadratischen
Gleichungen aber wäre die Matrix $V$ vierzeilig gewesen, und die
Bedingung, dass vier lineare homogene Gleichungen mit vier Unbekannten
durch ein bestimmtes Werthsystem derselben befriedigt werden können,
ist gerade durch das Verschwinden dieser Determinante $V$ ausgedrückt.

313. \textit{Berechnung von $p^2$ und $\pi^2$.} Unsere nächste Aufgabe
ist, $p^2$ als lineare Function quadratischer Covarianten
darzustellen. Es war
\[
p = (\beta\Delta)^2\beta_x = (\alpha\Delta)^2\alpha_x.
\]
Führen wir für $\Delta$ seinen Werth $(ab)^2a_xb_x$ ein, so wird:
\[
p = (ab)^2(a\beta)(b\beta)\beta_x.
\]
Multipliciren wir die beiden so erhaltenen Werthe für $p$, so kommt:
\[
p^2 = (\alpha\Delta)^2\alpha_x(ab)^2(a\beta)(b\beta)\beta_x.
\]
Hierin fordern die beiden quadratischen Klammerfactoren zur Benutzung
der Identität auf:
\[
(\alpha\Delta)^2(ab)^2 = \{(a\alpha)(b\Delta) - (b\alpha)(a\Delta)\}^2.
\]
Wegen der rechts auftretenden Factoren $(b\Delta)^2$, $(a\Delta)^2$
reducirt sich also der Werth von $p^2$ nach Einführung dieses Werthes
von
%CONT>>


% ===== pdf idx 351, printed 337.0 =====
% !PRINTED 337
\weberpage{337}
%<<CONT
$(a\Delta)^2(ab)^2$ auf:
\[
p^2 = -2(a\alpha)(b\alpha)(a\beta)(b\beta)(a\Delta)(b\Delta)\alpha_x\beta_x.
\]
Diesen Ausdruck wird man wegen der vier ersten Factoren mit Hilfe des Productsatzes umformen. Man erhält:
\[
p^2 = -(a\Delta)(b\Delta)\alpha_x\beta_x\left\{(a\alpha)^2(b\beta)^2+(a\beta)^2(b\alpha)^2-(ab)^2(\alpha\beta)^2\right\}.
\]
Die beiden ersten Glieder rechts gehen durch Vertauschung von $a$ mit $b$ in einander über, sind also einander gleich; daher:
\weq{6}{p^2 = (ab)^2(\alpha\beta)^2(a\Delta)(b\Delta)\alpha_x\beta_x - 2(a\alpha)^2(b\beta)^2(a\Delta)(b\Delta)\alpha_x\beta_x.}
Der erste Term rechts ist nichts anderes als das Product
\weq{7}{(ab)^2(a\Delta)(b\Delta)(\alpha\beta)^2\alpha_x\beta_x = A_{\Delta\Delta}\cdot \nabla.}
Um den zweiten Term umzuformen, benutzen wir wegen des Factors $(a\alpha)^2$ die Reihenentwicklung:
\weq{8}{(a\alpha)^2 a_y \alpha_x = \Theta_x\Theta_y + \frac{1}{2}J(yx).}
Denken wir uns in ihr $y$ durch $\Delta$ ersetzt und alsdann den Werth von $(a\alpha)^2(a\Delta)\alpha_x$ in den zweiten Term (6) eingetragen, so wird derselbe:
\[
-2(a\alpha)^2(a\Delta)\alpha_x\cdot(b\beta)^2(b\Delta)\beta_x = -2(b\beta)^2(b\Delta)\beta_x\left\{\Theta_x(\Theta\Delta)+\frac{1}{2}J\Delta_x\right\}.
\]
Transformiren wir die Factoren vor der Klammer nochmals durch dieselbe Reihenentwicklung (8), so kommt:
\weq{9}{\begin{aligned}
&-2(a\alpha)^2(a\Delta)\alpha_x\cdot(b\beta)^2(b\Delta)\beta_x\\
&= -2\Theta_x(\Theta\Delta)\cdot\left\{\Theta'_x(\Theta'\Delta)+\frac{1}{2}J\Delta_x\right\} - J\Delta_x\left\{\Theta_x(\Theta\Delta)+\frac{1}{2}J\Delta_x\right\}\\
&= -2(\Theta\Delta)\Theta_x(\Theta'\Delta)\Theta'_x - 2J\Delta_x(\Theta\Delta)\Theta_x - \frac{J^2}{2}\Delta\\
&= -2(\Theta\Delta)(\Theta'\Delta)\Theta_x\Theta'_x - 2J(\Theta,\Delta) - \frac{J^2}{2}\Delta.
\end{aligned}}
Das erste Glied rechts ist aber nach dem Productsatze:
\weq{10}{\begin{aligned}
2(\Theta\Delta)(\Theta'\Delta)\Theta_x\Theta'_x &= \left\{(\Theta\Delta)^2\Theta' + (\Theta'\Delta)^2\Theta - (\Theta\Theta')^2\Delta\right\}\\
&= 2A_{\Theta\Delta}\cdot\Theta - A_{\Theta\Theta}\cdot\Delta.
\end{aligned}}
Es wird also, wenn wir die in (7), (9) und (10) berechneten Werthe in (6) eintragen und überdies $-2(\Theta,\Delta)$ durch $(f,p)$ ersetzen [vergl. Nr. 310, (1)]
\weq{11}{p^2 = A_{\Delta\Delta}\nabla - 2A_{\Theta\Delta}\Theta + \left(A_{\Theta\Theta}-\frac{J^2}{2}\right)\Delta + J(f,p),}
und demgemäss, da $J$ bei Vertauschung von $a$ mit $\alpha$ das Zeichen wechselt:
\weq{12}{\pi^2 = A_{\nabla\nabla}\Delta - 2A_{\Theta\nabla}\Theta + \left(A_{\Theta\Theta}-\frac{J^2}{2}\right)\nabla - J(\varphi,\pi).}


% ===== pdf idx 352, printed 338.0 =====
% !PRINTED 338
\weberpage{338}
314. \textit{Exposition der folgenden Aufgaben.} Die vorausgegangenen Abschnitte dienen dazu, uns über einige Eigenschaften jener beiden linearen Covarianten $p$ und $\pi$ zu instruiren, die im ganzen System eine so hervorragende Rolle spielen. Wir können nun wiederum zu einer systematischen Methode zurückkehren, die uns über die Beziehungen der Formen des Systemes Aufschluss verschaffen soll. Indem wir diese Relationen zu ermitteln streben, werden wir uns gleichzeitig veranlasst sehen, gewisse Ueberschiebungen des Systemes durch ein geeignetes ihrer Glieder zu ersetzen. Wir beginnen damit, die Ueberschiebungen der einzelnen Covariantengruppen des Systemes der Reihe nach zu studiren, so weit sie für uns Interesse haben. Daran schliessen sich die Untersuchungen über die Relationen zwischen den Invarianten des Systemes, deren nothwendig zwei existiren müssen, da zwei cubische Formen nur $8-3=5$ unabhängige Constante besitzen, während doch 7 Invarianten im Systeme auftraten.

315. \textit{Die biquadratische Covariante $\vartheta$.} Die einzige Form vierter Ordnung, die im simultanen System von $f$ und $\varphi$ auftritt, ist die schiefe Covariante: $\vartheta = (f,\varphi) = (a\alpha)a_x^2\alpha_x^2$.

Von den Ueberschiebungen dieser Form über andere besprechen wir hier vier, die ein allgemeineres Interesse haben, nämlich die dritte Ueberschiebung über $f$ und $\varphi$, die erste über $p$ und $\pi$. Es ist:
\[
(\vartheta,f)^3 = \left((a\alpha)a_x^2\alpha_x^2,\, b_x^3\right)^3 = \frac{1}{2}(a\alpha)\left\{(ab)^2(\alpha b)\alpha_x + (\alpha b)^2(ab)a_x\right\}
\]
\[
= \frac{1}{2}(a\alpha)(ab)(\alpha b)\left\{(ab)\alpha_x + (\alpha b)a_x\right\}.
\]
Das letzte Glied $(\alpha b)a_x$ in der Klammer kann durch $\frac{1}{2}(ab)\alpha_x$ ersetzt werden, da die Identität besteht:
\[
(\alpha b)a_x + (a\alpha)b_x = (ab)\alpha_x
\]
und $(\alpha b)a_x$ hier mit $(a\alpha)b_x$ gleichbedeutend ist. Man erhält also:
\[
(\vartheta,f)^3 = \frac{3}{4}(a\alpha)(ab)^2(\alpha b)\alpha_x,
\]
oder, da $\Delta = (ab)^2a_xb_x$:
\weq{1}{(\vartheta,f)^3 = \frac{3}{4}(\Delta\alpha)(\alpha\Delta)\alpha_x = -\frac{3}{4}(\varphi,\Delta)^2 = -\frac{3}{4}p.}
Folglich auch, da $\vartheta$ durch Vertauschung von $a$ mit $\alpha$ sein Zeichen ändert:
\[
(\vartheta,\varphi)^3 = +\frac{3}{4}\pi.
\]
Zur Berechnung der ersten Ueberschiebung von $\vartheta$ über $p$ und $\pi$ gehen wir aus von der Fundamentalrelation Nr. 310, (5)
\[
a_x^2(a\pi) + \alpha_x^2(\alpha p) = 0.
\]


% ===== pdf idx 353, printed 339.0 =====
% !PRINTED 339
\weberpage{339}
Schieben wir diese Relation einmal über $f$, so kommt:
\weq{2}{0 = (ab)(a\pi)a_xb_x^2 + (\alpha b)(\alpha p)\alpha_xb_x^2.}
Das zweite Glied geht aus der Reihe
\[
(a\alpha)a_x^2\alpha_x\alpha_y = \vartheta_y - \frac{1}{2}\Theta(yx)
\]
hervor, wenn wir in ihr $a$ durch $b$ und $y$ durch $p$ ersetzen; es folgt dann
\weq{3}{(\alpha b)(\alpha p)a_xb_x^2 = -(\vartheta,p) + \frac{1}{2}\Theta p.}
Vertauschen wir im ersten Glied der Relation (2) $a$ mit $b$ und nehmen die halbe Summe, so erhalten wir:
\weq{4}{(ab)(a\pi)a_xb_x^2 = \frac{1}{2}(ab)a_xb_x\left\{(a\pi)b_x - (b\pi)a_x\right\} = \frac{1}{2}(ab)^2a_xb_x\cdot\pi.}
Tragen wir diese Werthe der beiden Glieder aus (3) und (4) in (2) ein, so erhalten wir:
\[
0 = \frac{\pi\Delta}{2} - (\vartheta,p) + \frac{1}{2}\Theta p,
\]
oder:
\weq{5}{(\vartheta,p) = \frac{1}{2}\left\{\Delta\pi + \Theta p\right\},}
und demgemäss:
\weq{6}{(\vartheta,\pi) = -\frac{1}{2}\left\{\nabla p + \Theta \pi\right\}.}
316. \textit{Die cubischen und quadratischen Covarianten.} Die Formen dritter Ordnung, welche das simultane System aufweist, wollen wir genau so beibehalten, wie sie der Faltungsprocess bei Aufstellung des Systemes geliefert hat. Es sind sechs an der Zahl:
\[
f,\ \varphi,\ Q,\ K,\ (f,\nabla),\ (\varphi,\Delta).
\]
Zu weitern Untersuchungen von Ueberschiebungen derselben über einander oder über andere Formen haben wir hier keine Veranlassung.

Von den sechs quadratischen Formen, die im Systeme auftreten, wollen wir drei durch andere ersetzen, nämlich die drei Formen $(\Delta,\nabla)$, $(Q,\varphi)^2$ und $(K,f)^2$. Die Gründe, die uns hiezu bewegen, sind rein praktischer Natur. \gesperrt{Clebsch} hatte dieselben ersetzt durch die drei Functionaldeterminanten der drei andern quadratischen Formen $\Delta$, $\nabla$, $\Theta$. Wir wollen an ihrer Stelle einführen: $(f,\pi)$, $(f,p)$, $(\varphi,\pi)$, also Ueberschiebungen über lineare Formen. Dass wir $(\Delta,\nabla)$ durch $(f,\pi)$ ersetzen können, geht aus Nr. 311, (3) hervor.

Berechnet man ferner die zweite Ueberschiebung von $Q$ über $\varphi$, so findet man Folgendes. Die zweite Polare von $Q$ ist [vgl. Nr. 146, (2)]:
\[
Q_xQ_y^2 = (a\Delta)a_y^2\Delta_x,
\]


% ===== pdf idx 354, printed 340.0 =====
% !PRINTED 340
\weberpage{340}
Ersetzen wir darin $y$ durch $\alpha$, und multipliciren mit $\alpha_x$, so kommt:
\[
(Q,\varphi)^2 = (a\Delta)(a\alpha)^2\Delta_x\alpha_x.
\]
Der Ausdruck rechts entsteht aber aus $(a\alpha)^2a_y\alpha_x$, wenn wir darin $y$ durch $\Delta$ ersetzen und mit $\Delta_x$ multipliciren. Wir erhalten [vergl. Nr. 313, (8)] daher:
\[
(Q,\varphi)^2 = (\Theta,\Delta) + \frac{J}{2}\Delta,
\]
also wegen Nr. 311, (1)
\weq{1}{(Q,\varphi)^2 = -\frac{1}{2}(f,p) + \frac{1}{2}J\Delta.}
Man darf also in der That $(Q,\varphi)^2$ durch $(f,p)$ ersetzen und demgemäss auch $(K,f)^2$ durch $(\varphi,\pi)$. Die sechs quadratischen Formen sind also:
\[
\Delta,\ \Theta,\ \nabla,\ (f,p),\ (\varphi,\pi),\ (f,\pi) = -(\varphi,p).
\]
317. \textit{Die sechs linearen Formen.} Von den sechs linearen Formen des Systemes behalten wir die beiden Formen $p$ und $\pi$ im Folgenden bei. Dagegen ersetzen wir die vier übrigen durch
\[
(\Delta,p),\ (\Delta,\pi),\ (\nabla,p),\ (\nabla,\pi).
\]
Dass dies gestattet ist, kann man leicht zeigen. Denn da die zweite Polare von $Q$, wie schon erwähnt, sich auf $(a\Delta)a_y^2\Delta_x$ reducirt, so wird
\[
(Q,\nabla)^2 = (a\Delta)(a\nabla)^2\Delta_x.
\]
Nun ist aber die lineare Form
\[
\pi_x = (a\nabla)^2a_x,
\]
also
\weq{1}{(Q,\nabla)^2 = \left((a\nabla)^2a_x,\, \Delta\right) = (\pi,\Delta)}
und ebenso
\weq{2}{(K,\Delta)^2 = \left((\alpha\Delta)^2\alpha_x,\, \nabla\right) = (p,\nabla).}
Man erhält ferner:
\[
(\varphi,\Delta^2)^3 = (\alpha\Delta)^2(\alpha\Delta_1)\Delta_{1x},
\]
oder, weil:
\[
(\alpha\Delta)^2\alpha_x = p,
\]
\weq{3}{(\varphi,\Delta^2)^3 = \left((\alpha\Delta)^2\alpha_x,\, \Delta_1\right) = (p,\Delta),}
und ebenso
\weq{4}{(f,\nabla^2)^3 = \left((a\nabla)^2a_x,\, \nabla_1\right) = (\pi,\nabla).}
Die sechs linearen Covarianten des Systemes sind also:
\[
p,\ \pi,\ (\Delta,p),\ (\Delta,\pi),\ (\nabla,p),\ (\nabla,\pi).
\]
318. \textit{Weitere lineare Covarianten.} Den eben besprochenen sechs Formen, insbesondere den beiden Ueberschiebungen von $f$ und $\varphi$ über
%CONT>>


% ===== pdf idx 355, printed 341.0 =====
